\capitulo{1}{Introducción}

En la última década, el interés por la enseñanza de la programación y la robótica ha crecido de manera significativa, con especial énfasis en la educación preuniversitaria, dirigida a \textbf{niños y adolescentes}. Esto se debe a la creciente necesidad de dotar a las nuevas generaciones de habilidades fundamentales en pensamiento computacional, resolución de problemas y lógica algorítmica, consideradas esenciales para su desarrollo en el mundo digital actual \cite{Paredes2022}. La alfabetización digital y el pensamiento computacional son habilidades transversales cruciales que promueven la resolución creativa de problemas y preparan a los estudiantes para los desafíos de la sociedad moderna \cite{Klass2023_BenefitsCodeKids, CodeWeek2020_EarlyCode}.

Entre las metodologías más empleadas en la enseñanza de la programación, la \textbf{programación basada en bloques} ha demostrado ser una herramienta pedagógica eficaz. Este paradigma visual facilita la comprensión de la lógica de programación y el flujo de control sin la complejidad de la sintaxis textual, reduciendo la barrera de entrada para estudiantes noveles y promoviendo un aprendizaje más accesible y lúdico \cite{Paredes2023, Munoz2015}. Su efectividad se ha demostrado en diversos contextos educativos, como la incorporación de robótica en el currículo STEM a través de herramientas de bloques como Scratch \cite{Fuchs2021_ScratchRobotics}. Experiencias en la introducción a la programación en la universidad han mostrado que el uso de estos lenguajes ha mejorado significativamente las tasas de aprobación en asignaturas fundamentales de la informática \cite{Munoz2015}.

Actualmente, dos de las plataformas más destacadas en este ámbito son \textbf{Scratch y Blockly}. Scratch, desarrollado por el MIT Media Lab, se ha consolidado como una de las herramientas educativas más utilizadas globalmente para la enseñanza de la programación, proporcionando un entorno altamente accesible y motivador \cite{Paredes2023}. Por otro lado, Blockly, desarrollado por Google, ofrece una biblioteca de código abierto que permite integrar programación por bloques en diversas plataformas y puede convertir los bloques visuales a lenguajes de programación textuales como JavaScript y Python, ofreciendo así una mayor versatilidad en la generación de código y aplicaciones \cite{Munoz2015}. Aunque \textbf{Scratch utiliza una implementación personalizada basada en los principios de Blockly}, existen diferencias fundamentales entre ambos sistemas, especialmente en términos de flexibilidad y aplicación práctica \cite{Campana2016}. La adopción y los avances en este campo se observan con gran desarrollo en diversas regiones de América Latina, donde la robótica educativa se ha posicionado como un medio eficaz para la enseñanza de habilidades STEM en docentes de primaria, evidenciando un compromiso temprano con estas metodologías \cite{Torres2021_LATAMSTEM}.

El \textbf{ aprendizaje basado en retos y la inclusión de elementos lúdicos } ha emergido como un enfoque clave en la enseñanza de la programación. Al introducir  \textit{desafíos con objetivos claros y un entorno interactivo}, se promueve una motivación intrínseca y un compromiso activo del estudiante en el proceso de aprendizaje, logrando así una experiencia más efectiva y significativa \cite{Campana2016}. Estudios recientes confirman que los proyectos de robótica aumentan significativamente la motivación y el engagement de los estudiantes, facilitando un aprendizaje más profundo y aplicado \cite{Dominguez2024_MotivationRobotics}. Además, \textit{los entornos interactivos con la programación (incluso con componentes de IA)}  demuestran resultados positivos en el rendimiento académico y la motivación en estudiantes de primaria y secundaria \cite{Alshahrani2024_GamifiedBlock}. Es en este cruce entre la rigorosa aplicación de los conocimientos técnicos y la creatividad del diseño, donde reside la oportunidad de desarrollar herramientas capaces de \textbf{vincular y emocionar a los estudiantes}, transformando el aprendizaje de la programación en una actividad lúdica y accesible que genere un impacto duradero. 

No obstante, a pesar de estos avances metodológicos y didácticos, existen desafíos significativos para la enseñanza de la programación y robótica, especialmente en entornos con acceso limitado a recursos educativos y tecnológicos. La falta de laboratorios físicos y la dependencia de materiales y robots costosos restringen la accesibilidad de estas valiosas herramientas, afectando principalmente a estudiantes en zonas rurales o con restricciones económicas \cite{Paredes2022, Lemus2006_JENUI}. La inviabilidad de replicar laboratorios con equipamiento sofisticado en contextos de bajo presupuesto y la dificultad de acceder a prácticas con componentes reales impulsan la búsqueda de alternativas flexibles y accesibles \cite{Paredes2022}. Esta problemática está directamente alineada con los \textbf{Objetivos de Desarrollo Sostenible (ODS)} de la Organización de las Naciones Unidas (ONU), en particular el \textbf{ODS 4 (Educación de calidad)}, que busca garantizar una educación inclusiva, equitativa y de calidad para niños y adolescentes \cite{ONU-ODS}. Se enfoca en que nadie se quede atrás en el acceso a oportunidades educativas, y el \textbf{ODS 10 (Reducción de las desigualdades)}, que promueve la reducción de las brechas sociales y económicas entre regiones y poblaciones, buscando asegurar un acceso equitativo al conocimiento y a las oportunidades educativas \cite{ONU-ODS}.

Frente a esta problemática, los \textbf{entornos de simulación virtual y los laboratorios virtuales de robótica} han demostrado ser una solución excepcionalmente eficaz para \textit{democratizar el acceso} a la enseñanza de la programación y robótica. Estas plataformas permiten a los estudiantes diseñar, construir y programar robots virtuales en un entorno digital, eliminando la necesidad de hardware físico costoso y ofreciendo un espacio de experimentación seguro y flexible. Permiten proporcionar una experiencia similar a la obtenida en un laboratorio de prácticas a través de sistemas de apoyo con gran realismo \cite{Simulador3DUnity}. El uso de \textbf{juegos serios} (serious games) y simulaciones interactivas se ha consolidado como una metodología prometedora para la enseñanza de la programación, al integrar desafíos divertidos y elementos educativos que mejoran la retención del conocimiento e incrementan significativamente la motivación y la participación activa de los estudiantes \cite{GeoBotsVR}.

 \textbf{El enfoque inicial de este proyecto consideraba la replicación completa de un sistema de programación puramente visual basado en bloques, similar a Scratch, dentro de Unity.} Se diseñó una arquitectura robusta para la gestión, conexión y ejecución de bloques, sentando unas bases técnicas sólidas. Sin embargo, durante el desarrollo, se constató que la implementación de un motor de layout gráfico completamente dinámico para las vistas de los bloques presentaba una complejidad técnica que desviaba el foco del objetivo principal: la lógica de programación y la simulación del robot. Ante este reto de ingeniería, se tomó una decisión estratégica: pivotar hacia un Producto Mínimo Viable (MVP) que, aun manteniendo la filosofía de simplicidad y la representación visual del programa, se centrase en una interacción textual. Este reenfoque permitió concentrar los esfuerzos en el diseño de un lenguaje específico, su intérprete y la interacción con el entorno 3D.



\textbf{Si bien el desarrollo de tecnologías en áreas como la simulación 3D o los lenguajes visuales es sólida, la \textit{integración holística y accesible} de estos elementos en un \textit{único sistema orientado a la robótica educativa específica de bajo coste} sigue presentando oportunidades significativas para el desarrollo y la implementación de soluciones innovadoras.}  La creación de un puente pedagógico entre los sistemas puramente visuales y la programación textual tradicional, mediante un lenguaje de comandos simple y un entorno 3D inmediato, representa una oportunidad significativa. Este proyecto busca aportar valor en este nicho, ofreciendo una herramienta que facilita la transición hacia la codificación escrita junto con la programación por bloques.


En este contexto, el presente \textbf{Trabajo Fin de Grado} se orienta al desarrollo de un \textbf{prototipo de laboratorio virtual de programación y robótica}. La plataforma está diseñada con el motor de videojuegos \textbf{Unity} utilizando \textbf{C\#}, con el objetivo de ofrecer una alternativa \textit{accesible, interactiva y didáctica} para la metodología \textbf{STEAM}. La herramienta introduce un lenguaje de comandos textual y un intérprete propio, diseñados para ser intuitivos y controlar un robot virtual en un entorno 3D. El programa construido por el usuario se representa visualmente como una lista secuencial de instrucciones, manteniendo la claridad de un sistema de bloques pero fomentando la práctica con la escritura de código. El entorno de simulación está inspirado en la funcionalidad de kits como LEGO WeDo 2.0 y está optimizado para su uso en dispositivos Android tipo tablet.


Este proyecto  plantea un desafío técnico centrado en el análisis sintáctico (parsing) y la interpretación de un lenguaje a medida. El sistema implementa un pipeline completo: desde la entrada de texto en una interfaz de usuario, su análisis y conversión a una estructura de datos intermedia (la lista de instrucciones), hasta la ejecución de dichas instrucciones en el motor 3D, gestionando el estado del robot y las interacciones con el entorno, como las colisiones.

La estructura de esta memoria detallará, en primer lugar, los objetivos específicos del prototipo. Posteriormente, se presentarán los conceptos teóricos que sustentan el desarrollo, incluyendo la arquitectura del intérprete de comandos y del motor de simulación. A continuación, se describirán las técnicas y herramientas empleadas. El capítulo central detallará los aspectos más relevantes del diseño e implementación del prototipo, seguido de una comparativa con trabajos relacionados. Finalmente, se expondrán las conclusiones obtenidas y las líneas de trabajo futuras para la evolución de este laboratorio virtual.